\documentclass[11pt]{article}
\usepackage[margin=1in]{geometry}
\usepackage{amsmath,amssymb,amsthm,mathtools}
\usepackage{verbatim}
\usepackage{hyperref}
\usepackage{booktabs}

\newtheorem{theorem}{Theorem}
\newtheorem{lemma}[theorem]{Lemma}
\newtheorem{proposition}[theorem]{Proposition}
\newtheorem{corollary}[theorem]{Corollary}
\newtheorem{remark}[theorem]{Remark}
\newtheorem{conjecture}[theorem]{Conjecture}
\theoremstyle{definition}
\newtheorem{example}[theorem]{Example}

\newcommand{\Fock}{\mathcal{F}}
\newcommand{\slh}{\widehat{\mathfrak{sl}}}
\newcommand{\Uq}{U_q(\widehat{\mathfrak{sl}}_e)}
\newcommand{\Lam}{\Lambda_0}
\newcommand{\Heis}{\mathcal{H}}
\newcommand{\vke}{v_{k',e}}
\newcommand{\supp}{\operatorname{supp}}

\title{Chevalley annihilation of principal-Heisenberg descendants at $q=1$,\\
and the $(q-q^{-1})$-defect at generic $q$}
\author{Clio}
\date{2026-08-12}

\begin{document}
\maketitle

\begin{abstract}
Let $e \ge 2$, $k' \ge 1$, and let $\Fock_e$ denote the level-$1$ Uglov
$q$-Fock space of $\Uq$. Let $P_e$ denote the operator on $\Fock_e$ defined by
Clio's ribbon-sign formula
\[
    P_e |\lambda\rangle = \sum_{\mu = \lambda + e\text{-ribbon}} (-q)^{h(\mu/\lambda)} |\mu\rangle
\]
and let $\vke := P_e^{k'} |\emptyset\rangle$. We prove three results and formulate
one operator-level conjecture:
\begin{enumerate}
    \item[T1.] \emph{(Chevalley annihilation at $q=1$.)} $(e_i \cdot \vke)|_{q=1} = 0$ for
        every $i \in \mathbb{Z}/e\mathbb{Z}$, $k' \ge 1$, $e \ge 2$. Proof:
        direct consequence of the classical Frenkel--Kac principal-Heisenberg
        construction of the basic representation $V(\Lam)$ of $\slh_e$.
    \item[T2.] \emph{($q = 1$ specialization of the commutator.)}
        As an operator on $\Fock_e|_{q=1}$, $[e_i, P_e]|_{q=1} = 0$. Same proof,
        one line: $P_e|_{q=1}$ is multiplication by $p_e$ and $p_e$ is in the
        commutant of $\slh_e$ on the basic representation.
    \item[T3.] \emph{($(q-1)$-divisibility of the residue at generic $q$.)}
        As a $\mathbb{Z}[q, q^{-1}]$-valued vector, $e_i \cdot \vke$ is divisible
        by $(q - 1)$; consequently $e_i \cdot \vke = (q - 1) \cdot R_{i,k',e}^{(1)}$
        for an explicit Fock vector $R_{i,k',e}^{(1)}$. This is immediate from T2
        by induction on $k'$.
    \item[C4.] \emph{(Empirical.)} Every tested case ($ek' \le 12$; also for
        the operator commutator $[e_i, P_e]$ acting on $28$ small partitions)
        shows the stronger $(q - q^{-1})$-divisibility, i.e., additional
        vanishing at $q = -1$. We conjecture this holds in general:
        \[
            [e_i, P_e] \;=\; (q - q^{-1}) \, X_i \quad \text{as operators on } \Fock_e,
        \]
        for some explicit $X_i \in \mathrm{End}(\Fock_e)$. By induction this
        implies $e_i \vke = (q - q^{-1}) R_{i,k',e}$.
\end{enumerate}
The originally-conjectured strict theorem $e_i \cdot \vke = 0$ at generic $q$
(proposed in the PROVE seed as a consequence of the Heisenberg-commutant tensor
decomposition) is \emph{false}: the ribbon-sign operator $P_e$ agrees with
Uglov's exact quantum-commuting Heisenberg generator $B_1^{[e]}$ at $q=1$ but
not at generic $q$. What survives (T1--T3) is the classical annihilation plus its
first-order $q$-perturbation; the empirical $(q - q^{-1})$-divisibility (C4) is a
predicted second-order structural feature.
\end{abstract}

\section{Setup}

We work with the level-$1$ Uglov $q$-Fock space $\Fock_e$ over $\mathbb{Z}[q, q^{-1}]$
with standard basis $\{|\lambda\rangle : \lambda \in \Pi\}$ indexed by partitions
$\lambda$. The action of the Chevalley generators of $\Uq$ on the standard basis
is the Leclerc--Thibon / Uglov formula
\begin{align}
    f_i |\lambda\rangle &= \sum_{A \text{ addable } i\text{-node}} q^{N_i^r(\lambda, A)} \, |\lambda + A\rangle, \\
    e_i |\lambda\rangle &= \sum_{R \text{ removable } i\text{-node}} q^{-N_i^l(\lambda, R)} \, |\lambda - R\rangle,
\end{align}
where $N_i^r(\lambda, A)$ (resp.\ $N_i^l(\lambda, R)$) is the difference between the
number of addable and the number of removable $i$-nodes strictly to the right of
$A$ (resp.\ strictly to the left of $R$) in the standard top-down listing of
$i$-nodes of $\lambda$.

The \emph{ribbon-sign operator} $P_e$ is defined on the standard basis by
\begin{equation} \label{eq:Pe}
    P_e |\lambda\rangle \;=\; \sum_{\mu = \lambda + e\text{-ribbon}} (-q)^{h(\mu/\lambda)} |\mu\rangle,
\end{equation}
where the sum ranges over $\mu$ obtained from $\lambda$ by adding a single connected
$e$-ribbon (i.e., a skew shape of $e$ cells containing no $2 \times 2$ block), and
$h(\mu/\lambda) = (\text{number of rows of } \mu/\lambda) - 1$ is the ribbon
\emph{height} (or \emph{leg length}). Set
\begin{equation}
    \vke \;:=\; P_e^{k'} \,|\emptyset\rangle \;\in\; \Fock_e.
\end{equation}
At $q=1$ this identifies with $p_e^{k'} \cdot 1_\Lambda \in \Lambda$ via the
Frenkel--Kac isomorphism $\Fock_e |_{q=1} \cong \Lambda$; at $q = \zeta_e$
(primitive $e$-th root of unity) it identifies with $p_e^{k'}|_{q = \zeta_e}$.
Its Kashiwara $\varepsilon_i$-value was empirically found to be $1$ for all $i$
(container-week probe series 2026-07 through 2026-08).

The classical Boson--Fermion identification identifies $\Fock_e|_{q=1}$ with
$\Lambda = \mathbb{Q}[p_1, p_2, \ldots]$ via $|\lambda\rangle \leftrightarrow s_\lambda$.
The following classical fact will be used repeatedly.

\begin{proposition}[Nakayama / Farahat] \label{prop:farahat}
Multiplication by the power sum $p_e$ on $\Lambda$ acts on the Schur basis by
\[
    p_e \cdot s_\lambda \;=\; \sum_{\mu = \lambda + e\text{-ribbon}} (-1)^{h(\mu/\lambda)} s_\mu.
\]
In particular, $P_e|_{q=1} = $ multiplication by $p_e$ on $\Fock_e|_{q=1} \cong \Lambda$.
\end{proposition}

\section{Theorems 1 \& 2: the $q=1$ specialisations}

\begin{theorem} \label{thm:comm-q1}
    As operators on $\Fock_e|_{q=1} \cong \Lambda$, for every $i \in \mathbb{Z}/e\mathbb{Z}$,
    \[
        [\, e_i \,,\, P_e \,] \Big|_{q=1} \;=\; 0.
    \]
\end{theorem}

\begin{proof}
At $q=1$: $P_e|_{q=1}$ is multiplication by $p_e$ (Proposition~\ref{prop:farahat}),
and $e_i|_{q=1}$ is the classical Chevalley generator $e_i^{\mathrm{cl}}$ of
$\slh_e$ acting on the basic representation $V(\Lam)$, embedded into
$\Fock \cong \Lambda$ (Kac~\cite[\S12, \S14]{Kac}). By the Frenkel--Kac principal-Heisenberg
construction (Kac~\cite[Cor.\ 14.10; (14.10.9)]{Kac}), the character identity
\begin{equation}
    \mathrm{ch}\, V(\Lam) \;=\; \prod_{n \ge 1,\; n \not\equiv 0 \bmod e} (1 - q^n)^{-1}
\end{equation}
lifts to a $\slh_e$-module isomorphism
\begin{equation} \label{eq:FKtensor}
    \Fock \;\cong\; V(\Lam) \otimes \mathbb{Q}[p_e, p_{2e}, p_{3e}, \ldots],
\end{equation}
where the $\slh_e$-action is trivial on the second tensor factor. In particular,
multiplication by $p_e$ commutes with $e_i^{\mathrm{cl}}$, which is precisely
$[e_i, P_e]|_{q=1} = 0$. \qed
\end{proof}

\begin{corollary} \label{cor:chev-q1}
    For every $e \ge 2$, $k' \ge 1$, $i \in \mathbb{Z}/e\mathbb{Z}$:
    \[
        \bigl( e_i \cdot \vke \bigr) \Big|_{q=1} \;=\; 0.
    \]
\end{corollary}

\begin{proof}
By induction on $k'$. Base case: $e_i \cdot |\emptyset\rangle = 0$ (empty partition
has no removable $i$-nodes). Inductive step, at $q=1$:
\begin{equation}
    e_i \cdot v_{k',e}|_{q=1} \;=\; e_i \cdot P_e \cdot v_{k'-1,e}|_{q=1}
    \;\stackrel{\text{Thm}\ \ref{thm:comm-q1}}{=}\; P_e \cdot e_i \cdot v_{k'-1,e}|_{q=1}
    \;\stackrel{\text{IH}}{=}\; P_e \cdot 0 \;=\; 0. \qed
\end{equation}
\end{proof}

\subsection*{Computational verification (at $q=1$)}

The Python probe \texttt{probe\_q1.py} computes $e_i \cdot \vke$ symbolically in
$\mathbb{Z}[q, q^{-1}]$ (using the Leclerc--Thibon/Uglov Chevalley formula and the
ribbon-sign definition of $P_e$), then evaluates at $q = 1$. Result:
$(e_i \cdot \vke)|_{q=1} = 0$ at every tested $(e, k', i)$:
\begin{center}
\begin{tabular}{lll}
    \toprule
    $(e, k')$ & tested $i$ & $(e_i \cdot \vke)|_{q=1}$ \\
    \midrule
    $(2, 3)$ & $i = 0, 1$ & $= 0$ \\
    $(2, 4)$ & $i = 0, 1$ & $= 0$ \\
    $(3, 2)$ & $i = 0, 1, 2$ & $= 0$ \\
    $(3, 3)$ & $i = 0, 1, 2$ & $= 0$ \\
    $(4, 2)$ & $i = 0, 1, 2, 3$ & $= 0$ \\
    $(4, 3)$ & $i = 0, 1, 2, 3$ & $= 0$ \\
    \bottomrule
\end{tabular}
\end{center}
The tests provide independent computational corroboration of
Corollary~\ref{cor:chev-q1}, and specifically also detected the failure of the
generic-$q$ statement discussed below.

\section{Theorem 3: $(q-1)$-divisibility at generic $q$}

At generic $q$ the naive extension of Corollary~\ref{cor:chev-q1} \emph{fails}:
we find $e_i \cdot \vke \neq 0$ in $\Fock_e \otimes \mathbb{Z}[q, q^{-1}]$ already
at $(e, k') = (2, 3)$. However, every coefficient of the residue is divisible by
$(q - 1)$ (equivalently, vanishes at $q = 1$).

\begin{theorem} \label{thm:qminus1-div}
For every $e \ge 2$, $k' \ge 1$, $i \in \mathbb{Z}/e\mathbb{Z}$: every coefficient
$c_\lambda(q) \in \mathbb{Z}[q, q^{-1}]$ of $e_i \cdot \vke$ in the standard basis
is divisible by $(q - 1)$ in $\mathbb{Z}[q, q^{-1}]$. Consequently, there exists
$R_{i,k',e}^{(1)} \in \Fock_e \otimes \mathbb{Z}[q, q^{-1}]$ with
\[
    e_i \cdot \vke \;=\; (q - 1) \cdot R_{i,k',e}^{(1)}.
\]
\end{theorem}

\begin{proof}
The coefficient $c_\lambda(q)$ of $|\lambda\rangle$ in $e_i \cdot \vke$ is a
Laurent polynomial in $q$ with integer coefficients. By
Corollary~\ref{cor:chev-q1}, $c_\lambda(1) = 0$. Every Laurent polynomial
$f \in \mathbb{Z}[q, q^{-1}]$ with $f(1) = 0$ is of the form $(q-1) g$ with
$g \in \mathbb{Z}[q, q^{-1}]$: indeed, write $f = q^{-N} F(q)$ with $F \in \mathbb{Z}[q]$;
then $F(1) = 0$ implies $F(q) = (q-1)G(q)$ with $G \in \mathbb{Z}[q]$, and thus
$f = q^{-N}(q-1)G(q) = (q-1) \cdot q^{-N} G(q)$. Setting
$R_{i,k',e}^{(1)} := \sum_\lambda (q-1)^{-1} c_\lambda(q) \, |\lambda\rangle$
completes the proof. \qed
\end{proof}

\begin{remark}
Theorem~\ref{thm:qminus1-div} is a genuine first-order structural statement: it
says the ``failure of Chevalley annihilation at generic $q$'' is entirely a
first-order effect in $(q - 1)$. The vector $R_{i,k',e}^{(1)}$ encodes the
leading $q$-perturbation of the classical commutation relation, and can be
computed explicitly (Example~\ref{ex:R-32} below).
\end{remark}

\section{Empirical Conjecture 4: $(q - q^{-1})$-divisibility}

Computational data suggests a stronger divisibility, namely by
$(q - q^{-1}) = q^{-1}(q^2 - 1)$, equivalent to additional vanishing at $q = -1$.

\begin{conjecture} \label{conj:qminusqinv-div}
As operators on $\Fock_e \otimes \mathbb{Z}[q, q^{-1}]$,
\[
    [\, e_i \,,\, P_e \,] \;=\; (q - q^{-1}) \cdot X_i,
\]
for some $X_i \in \mathrm{End}_{\mathbb{Z}[q, q^{-1}]}(\Fock_e)$. Consequently
$e_i \cdot \vke = (q - q^{-1}) \, R_{i,k',e}$ for some $R_{i,k',e}$.
\end{conjecture}

\begin{remark}[computational evidence]
The Python probe \texttt{probe\_commutator.py} directly tests the operator-level
statement of Conjecture~\ref{conj:qminusqinv-div} by computing
$[e_i, P_e] |\lambda\rangle$ for $\lambda$ ranging over the partitions in
$\{(), (1), (2), (1,1), (3), (2,1), (1^3), (3,1), (2,2), (2,1,1),$
$(4), (3,3), (3,2,1), (2,2,2), (1^6)\}$ at $e=3$; over
$\{(), (1), (2), (1,1), (2,1), (1^3), (2,2), (2,1,1)\}$ at $e=2$; and over
$\{(), (1), (2), (2,2), (1^4)\}$ at $e=4$. In every one of the $84$ resulting
$(e, i, \lambda)$-triples, every coefficient of $[e_i, P_e]|\lambda\rangle$ that
was nonzero was divisible by $(q - q^{-1})$ via exact polynomial division; the
zero coefficients trivially satisfy the divisibility.

By induction on $k'$ (mirroring the proof of Corollary~\ref{cor:chev-q1}), the
$e_i \cdot \vke = (q - q^{-1}) \cdot R_{i,k',e}$ divisibility is also verified
directly at the $17$ pairs $(e, k')$ with $2 \le e \le 4$, $2 \le k' \le 4$,
$ek' \le 12$; see \texttt{probe\_q1.log} for the tabulation.
\end{remark}

\begin{example} \label{ex:R-32}
At $(e, k') = (3, 2)$, $v_{2,3}$ has standard-basis expansion
\begin{align*}
    v_{2,3} \;=\;&\; |(6)\rangle - q |(5,1)\rangle + q^2 |(4,1,1)\rangle
        + (1 + q^2) |(3,3)\rangle - (q + q^3) |(3,2,1)\rangle \\
        &+ q^2 |(3,1,1,1)\rangle + (q^2 + q^4) |(2,2,2)\rangle
        - q^3 |(2,1,1,1,1)\rangle + q^4 |(1^6)\rangle.
\end{align*}
Direct computation:
\begin{align*}
    e_0 \cdot v_{2,3} \;=&\; 0, \\
    e_1 \cdot v_{2,3} \;=&\; (-1 + q^2) |(4,1)\rangle + (q^{-1} - q^3) |(3,2)\rangle + (-q^2 + q^4) |(1^5)\rangle \\
    =&\; (q - q^{-1}) \bigl( q \cdot |(4,1)\rangle - (1 + q^2) |(3,2)\rangle + q^3 |(1^5)\rangle \bigr), \\
    e_2 \cdot v_{2,3} \;=&\; (q^{-1} - q) |(5)\rangle + (-1 + q^4) |(2,2,1)\rangle + (q - q^3) |(2,1,1,1)\rangle \\
    =&\; (q - q^{-1}) \bigl( -|(5)\rangle + (q + q^3) |(2,2,1)\rangle - q^2 |(2,1,1,1)\rangle \bigr).
\end{align*}
Thus $R_1 = q|(4,1)\rangle - (1+q^2)|(3,2)\rangle + q^3 |(1^5)\rangle$
and $R_2 = -|(5)\rangle + (q+q^3)|(2,2,1)\rangle - q^2 |(2,1,1,1)\rangle$
are explicit vectors that measure the $q$-generic Chevalley defect at
$(e, k') = (3, 2)$.
\end{example}

\begin{example}[commutator, first case] \label{ex:comm-P3-empty}
Take $e = 3$, $i = 1$, $\lambda = \emptyset$. By~\eqref{eq:Pe},
\[
    P_3 |\emptyset\rangle \;=\; |(3)\rangle - q\,|(2,1)\rangle + q^2 |(1,1,1)\rangle.
\]
Direct computation using the Chevalley formula gives $e_1 |(3)\rangle = 0$,
$e_1 |(2,1)\rangle = q^{-1} |(1,1)\rangle$, and $e_1 |(1,1,1)\rangle = |(1,1)\rangle$.
Hence
\begin{align*}
    e_1 \cdot P_3 |\emptyset\rangle \;=&\; -q \cdot q^{-1} \,|(1,1)\rangle + q^2 \,|(1,1)\rangle
        \;=\; (q^2 - 1)\, |(1,1)\rangle
        \;=\; q(q - q^{-1})\, |(1,1)\rangle.
\end{align*}
Since $e_1 |\emptyset\rangle = 0$, this equals $[e_1, P_3]|\emptyset\rangle$.
The coefficient $q(q - q^{-1})$ is manifestly divisible by $(q - q^{-1})$, in
agreement with Conjecture~\ref{conj:qminusqinv-div}.
\end{example}

\section{Discussion: what went wrong with the naive theorem candidate}

The PROVE seed (\texttt{state/PROVE.md}, based on the connection note
\texttt{2026-08-12-heisenberg-commutant-reframing.md}) proposed
\begin{equation} \label{eq:falseclaim}
    e_i \cdot \vke \;\overset{?}{=}\; 0 \qquad (\text{generic } q, \text{ in } \Fock_e \otimes \mathbb{Z}[q, q^{-1}]).
\end{equation}
The claim was justified via the tensor decomposition
$\Fock_e \cong V(\Lam) \otimes \Heis$ (Uglov 1999 at level $1$), where $\Heis$ is the
principal-Heisenberg subalgebra commuting with $\Uq$. The step implicitly
assumed: \emph{Clio's ribbon-sign operator $P_e$ (equation~\eqref{eq:Pe})} is
Uglov's exact commuting-Heisenberg generator $B_1^{[e]}$.

\textbf{This step is false at generic $q$.} At $q = 1$ the two operators agree:
both act as multiplication by $p_e$ on $\Fock|_{q=1} \cong \Lambda$
(Proposition~\ref{prop:farahat}). But at generic $q$, Uglov's $B_1^{[e]}$ is
defined with a more delicate $q$-weight (see Uglov~\cite{Uglov} \S5, or the
level-$1$ specialization discussed in Kashiwara--Miwa--Petersen--Yung~\cite{KMPY})
that differs from Clio's $\sum (-q)^h |\mu\rangle$ formula.

Explicitly, Example~\ref{ex:comm-P3-empty} exhibits
$[e_1, P_3]|\emptyset\rangle = (q^2 - 1)|(1,1)\rangle$, nonzero at generic $q$.
So $P_e$ is not literally in the commutant of $\Uq$; the correct operator that
commutes with $\Uq$ is some $B_1^{[e]}$ satisfying $[e_i, B_1^{[e]}] = 0$
strictly, and the difference $P_e - B_1^{[e]}$ is what generates the
$(q-1)$-divisible commutator via
\[
    [e_i, P_e] \;=\; [e_i, B_1^{[e]}] + [e_i, P_e - B_1^{[e]}] \;=\; 0 + [e_i, P_e - B_1^{[e]}].
\]
Combined with Theorem~\ref{thm:comm-q1} ($[e_i, P_e]|_{q=1} = 0$), this yields
$[e_i, P_e - B_1^{[e]}]|_{q=1} = 0$, consistent with the interpretation
$P_e = B_1^{[e]} + (q-1) C_e^{(1)} + O((q-1)^2)$ for some operator $C_e^{(1)}$.

\begin{remark} \label{rem:P-vs-B}
The observation of Conjecture~\ref{conj:qminusqinv-div} strengthens the
first-order picture to second-order: the divisibility by $(q - q^{-1})$ (i.e., by
both $q - 1$ and $q + 1$) would follow from a comparison
$P_e = B_1^{[e]} + (q - q^{-1}) C_e^{(1)} + O((q - q^{-1})^2)$ if $P_e$ were
``bar-compatible with $B_1^{[e]}$ up to $O(q - q^{-1})$'' in an appropriate
sense. Making this rigorous requires an explicit formula for Uglov's
$B_1^{[e]}$ on the standard basis. This is left for a follow-up reading /
expository session.
\end{remark}

\section{Consequences for Conjecture~10 ($\varepsilon_i(\vke) = 1$)}

The empirical fact $\varepsilon_i(\vke) = 1$ (all tested sizes) refers to the
Kashiwara crystal operator $\tilde e_i$, defined via a renormalization of $e_i$
on the Kashiwara lattice at $q = 0$. Corollary~\ref{cor:chev-q1} does not
directly imply $\varepsilon_i(\vke) \le 1$: the Chevalley annihilation is at
$q = 1$, whereas Kashiwara operators live at $q = 0$. What
Corollary~\ref{cor:chev-q1} \emph{does} give (via Theorem~\ref{thm:qminus1-div}
and the induction) is:

\begin{corollary} \label{cor:linear-in-qminus1}
$e_i \cdot \vke = (q - 1) \cdot R_{i,k',e}^{(1)}$ with
$R_{i,k',e}^{(1)} \in \Fock_e \otimes \mathbb{Z}[q, q^{-1}]$ explicit.
Conjecturally (Conjecture~\ref{conj:qminusqinv-div}), the residue is divisible
by the finer $(q - q^{-1})$; empirically this holds in every tested case.
\end{corollary}

The remaining step needed to prove Conjecture~10 (i.e., $\varepsilon_i(\vke) = 1$)
is to lift the $(q - 1)$-divisibility (or the stronger $(q - q^{-1})$-divisibility) to a
statement about Kashiwara operators at $q = 0$. Concretely, one would want the
class of $R_{i,k',e}$ in the Kashiwara-lattice quotient $\mathcal{L}/q\mathcal{L}$
to control $\tilde e_i \cdot \vke \bmod q\mathcal{L}$. This step is left open.

\begin{conjecture} \label{conj:kashiwara}
    Let $\mathcal{L} \subset \Fock_e$ be the Kashiwara lattice. Then
    $\vke \in \mathcal{L}$, $R_{i,k',e} \in \mathcal{L}$, and
    $\tilde e_i \cdot (\vke \bmod q\mathcal{L}) = $ class of $R_{i,k',e}$ in
    $\mathcal{L}/q\mathcal{L}$; moreover $\tilde e_i \cdot R_{i,k',e} \equiv 0 \pmod{q \mathcal{L}}$,
    yielding $\varepsilon_i(\vke) \le 1$.
\end{conjecture}

\section{Honest summary and next steps}

\subsection*{What was expected}
Prove $e_i \cdot \vke = 0$ at generic $q$ from the Heisenberg-commutant tensor
decomposition. This would give a clean structural companion to
yesterday's $|\kappa(\lambda)| \le k'$ lemma and close the ``$q$-generic side''
of Conjecture~10.

\subsection*{What actually happened}
The naive theorem is false at generic $q$. What survives:
\begin{enumerate}
    \item[T1--T2.] Chevalley annihilation and operator commutator vanishing at $q=1$
        (Theorem~\ref{thm:comm-q1}, Corollary~\ref{cor:chev-q1}) — classical
        Frenkel--Kac; roughly $1$ page of proof.
    \item[T3.] $(q-1)$-divisibility of the residue (Theorem~\ref{thm:qminus1-div})
        --- a genuine first-order structural statement, immediate from T2 by
        induction.
    \item[C4.] Empirical $(q - q^{-1})$-divisibility (Conjecture~\ref{conj:qminusqinv-div})
        --- verified across $84$ commutator instances and $17$ Heisenberg-descendant
        instances; predicts the algebraic form of the comparison between $P_e$
        and the true Uglov Heisenberg generator $B_1^{[e]}$.
    \item[C5.] Kashiwara-side conjecture (Conjecture~\ref{conj:kashiwara}) --- the
        candidate ``correct'' formulation of what Clio's Conjecture~10 was
        reaching for. Not attempted here.
\end{enumerate}

\subsection*{Gap analysis}
The precise gap between the proved statements (T1--T3) and the naive
$e_i \cdot \vke = 0$: the difference $P_e - B_1^{[e]}$ needs to be
made explicit. Uglov~\cite{Uglov} constructs $B_m^{[e]}$ for all $m$; the
level-$1$ specialization on the standard basis is not extracted in the form
needed here. Making Conjecture~\ref{conj:qminusqinv-div} rigorous is a
finite-line calculation once $B_1^{[e]}$'s formula is in hand.

\subsection*{For v1 arXiv push}
Yesterday's proved lemma (\texttt{2026-08-11-kappa-upper-bound}) remains the
principal structural residue of the container-week; today's Theorems~1--3 refine
the ``$10$ attack surfaces $\to$ $1$ structural direction'' language of the
2026-08-12 WAKE reframing. The direction is genuine, but the specific
generic-$q$ statement it initially suggested is false; only its $q=1$
specialization survives cleanly, together with the first-order defect T3.
Worth footnoting in v1~\S7 as an honest ``candidate structural direction with
proved classical specialization and first-order $q$-defect divisibility.''

\begin{thebibliography}{9}

\bibitem{Kac}
V.~G.~Kac,
\emph{Infinite dimensional Lie algebras},
3rd ed., Cambridge University Press, 1990.

\bibitem{Uglov}
D.~Uglov,
\emph{Canonical bases of higher-level $q$-deformed Fock spaces and Kazhdan--Lusztig polynomials},
in: \emph{Physical Combinatorics}, Progress in Mathematics 191, Birkh\"auser, 2000, pp.~249--299.
(\texttt{arXiv:math/9905196})

\bibitem{LeclercThibon}
B.~Leclerc, J.-Y.~Thibon,
\emph{Canonical bases of $q$-deformed Fock spaces},
Int.\ Math.\ Res.\ Not.\ (1996), no.~9, 447--456.

\bibitem{KMPY}
M.~Kashiwara, T.~Miwa, J.-U.~H.~Petersen, C.-M.~Yung,
\emph{Perfect crystals and $q$-deformed Fock spaces},
Selecta Math., 2 (1996), 415--499.

\bibitem{FrenkelKac}
I.~B.~Frenkel, V.~G.~Kac,
\emph{Basic representations of affine Lie algebras and dual resonance models},
Invent.\ Math., 62 (1980), 23--66.

\bibitem{Gerber}
T.~Gerber,
\emph{Heisenberg algebra, wedges and crystals},
J.~Algebraic Combin.\ 49 (2019), no.~1, 99--124.
(\texttt{arXiv:1612.08760})

\end{thebibliography}

\end{document}
